\section{Formulación de la pregunta de investigación}
\label{sec:pregunta}

\subsection{Definición de la pregunta orientadora}

A partir del vacío de conocimiento identificado en la literatura técnica
(sección~\ref{sec:estado-arte}) y de la delimitación del problema analítico de
datos en el chancado primario, se formula la siguiente pregunta de
investigación:

\begin{quote}
\emph{¿De qué manera el diseño y la validación mediante simulación estocástica
de un módulo de software analítico basado en el algoritmo del Filtro de Kalman
Lineal permite mitigar el ruido instrumental y reconstruir de forma confiable el
estado dinámico interno de volumen de carga en un chancador primario
parcialmente observable?}
\end{quote}

\subsection{Alcance e integración de la pregunta}

La pregunta actúa como eje articulador central del anteproyecto y su formulación
delimita el estudio en tres sentidos simultáneos.

En cuanto al \textbf{objeto}, restringe el alcance exclusivamente a la capa de
software de carácter matemático y probabilístico, descartando de manera
explícita el desarrollo de interfaces comerciales, el despliegue de hardware y
la integración con protocolos industriales operativos.

En cuanto al \textbf{método}, el verbo «validar mediante simulación estocástica»
compromete un diseño experimental cuantitativo y no un desarrollo puramente
descriptivo: exige verificar cuantitativamente la confiabilidad de la
reconstrucción mediante métricas estadísticas explícitas dentro de un entorno
informático controlado.

En cuanto al \textbf{criterio de respuesta}, la expresión «de forma confiable»
no queda librada a juicio cualitativo. Se operacionaliza mediante la reducción
del error cuadrático medio respecto de una línea base declarada —la lectura
SCADA cruda— y mediante una prueba de hipótesis formal que descarte que dicha
reducción sea atribuible al azar estocástico.

De este modo, la pregunta vincula directamente la brecha de incertidumbre
identificada en el problema con el diseño metodológico propuesto y con los
criterios bajo los cuales los resultados serán evaluados, asegurando que la
respuesta sea verificable y no meramente argumentativa.
